\chapter{אימות הרעיון ומציאת בעיה אמיתית}

הטעות הנפוצה ביותר בתחילת הדרך היא להתאהב בפתרון במקום בבעיה. רעיונות טובים
לסטארט-אפ נובעים בדרך כלל מבעיה אמיתית שהיזם חווה בעצמו, ולא ממאמץ מכוון
``להמציא רעיון''~\cite{graham2005ideas}. בעיה שכדאי לפתור היא כזו שהיא
\emph{תכופה}, \emph{כואבת} ושהקהל מוכן לשלם כדי להיפטר ממנה.

\section{גילוי לקוחות לפני כתיבת קוד}

מתודולוגיית פיתוח הלקוחות \term{(Customer Development)} גורסת שיש לצאת מחוץ
לבניין ולדבר עם לקוחות פוטנציאליים \emph{לפני} שבונים מוצר~\cite{blank2013startup}.
מטרת השיחות אינה למכור, אלא ללמוד: מהי הבעיה, כיצד הלקוח פותר אותה היום, וכמה
היא עולה לו. שאלות טובות עוסקות בעבר ובהווה (``מתי נתקלת בבעיה לאחרונה?'') ולא
בעתיד ההיפותטי (``האם היית קונה?'').

\section{ניסוח השערות בדיקות}

רעיון הוא אוסף של השערות \term{(hypotheses)} שטרם נבדקו: מי הלקוח, מהי הבעיה,
מה הפתרון, ומדוע ישלמו עליו. עבודה נכונה ממירה כל השערה לניסוי זול ומהיר. אם
ההשערה החשובה ביותר — שקיים ביקוש — אינה מתאמתת, עדיף לגלות זאת תוך שבוע ובעלות
זניחה, ולא לאחר שנה של פיתוח. זהו לב ההיגיון הרזה: \emph{להקטין בזבוז} על ידי
למידה מהירה.
